% ===================================================================
%  CADRO N.º 7 — A aldea no inverno. — A casa labrega na primavera.
%  Traduction 2026 d'apres texto/io, controlee sur texto/fr, texto/es
%  et texto/pt. Consigne : voir texto/gl/10-cadro-01.tex.
%
%  OUVERTURE DE LA DEUXIEME SERIE : « SEGUNDA SERIE ».
%
%  LE TABLEAU DE LA FERME EST CELUI OU LE GALICIEN EST LE PLUS
%  LUI-MEME DE TOUTE LA COLONNE, et cela s'explique par la meme regle
%  que partout : la ferme s'apprend dehors, sans maitre, et de
%  personne qui vienne d'ailleurs. « Xugo » le joug, « eira » l'aire,
%  « sacho » la houe, « fouce » la faux, « angazo » le rastel,
%  « cortello » l'etable, « pombal » le colombier, « galiñeiro » le
%  poulailler, « esterco » le fumier : rien de tout cela n'est venu
%  d'une ecole ni d'un livre.
%
%  ET C'EST ICI QUE LE GALICIEN CESSE D'ETRE ENTRE SES DEUX VOISINES
%  POUR ETRE A COTE D'ELLES. Les mots de la ferme ne sont ni ceux de
%  Madrid ni ceux de Lisbonne : « leiteira » n'est pas « lechera »,
%  « cortello » n'est ni « establo » ni « estábulo », « rula » n'est
%  ni « paloma » ni « pomba ». Sur les six tableaux deja faits, la
%  distance a l'espagnol se creuse a mesure qu'on s'eloigne de
%  l'ecole : maximale au corps et a la ferme, minimale au chantier et
%  a la maladie.
% ===================================================================

%%K t07-noto-1 noto
\VUnotes{143.2mm}{%
(*) Véxanse os detalles da comida nos cadros 6 e 13.%
}

%%K t07-apar-1 apar
\VUsaut{4.0mm}
\VUcentre{13.2pt}{90}{SEGUNDA SERIE}
\VUsaut{3.0mm}
\VUfilet{16mm}
\VUsaut{5.6mm}
\VUtitre{15.0pt}{60}{CADRO N.\textsuperscript{o} 7}
\VUsaut{3.0mm}

%%K t07-tit-1 sub
\VUcentre{12.0pt}{0}{\textit{Primeira escena.}}
\VUsaut{3.0mm}

%%K t07-tit-2 sub
\VUpk{10.2pt}{40}{A aldea no inverno.}
\VUsaut{4.0mm}

%%K t07-c1-01-1 p
1. --- Durante as vacacións de Aninovo acompañei o meu pai a Vilanova, unha aldeíña
onde tiña que tratar algúns negocios de comercio.

%%K t07-c1-02-1 p
2. --- Collemos a \VUgras{dilixencia} \textsuperscript{(11)} que anda entre a vila e
ese burgo; pero, como facía moito frío, aquela viaxe nun vehículo pouco cómodo non foi
moi agradable.

%%K t07-c1-03-1 p
3. --- Por iso sentimos verdadeiro pracer cando vimos o \VUgras{condutor} parar o seu
carro na praza, diante da \VUgras{pousada} \textsuperscript{(2)} do \textit{Oso
Branco}. O \VUgras{pousadeiro} \textsuperscript{(4)} estábanos a agardar na soleira
da súa casa, fumando tranquilamente a súa \VUgras{pipa.}

%%K t07-c1-03-2 p
Convidounos a entrar e non me fixen de rogar para me instalar diante do lume de
sarmentos que crepitaba no lar. Entón burlei ben do agudo \VUgras{vento do norte} que
facía renxer o vello \VUgras{letreiro} \textsuperscript{(3)} e levaba en remuíños
negros \VUgras{o fume} \textsuperscript{(1)} que saía da cheminea.

%%K t07-c1-04-1 p
4. --- Era a hora do xantar. Sen agardar máis, comemos ben a comida sinxela pero
saborosa que nos serviron (*).

%%K t07-tit-3 sub
\VUpk{10.2pt}{40}{Algunhas visitas.}
\VUsaut{3.0mm}

%%K t07-c1-05-1 p
5. --- Despois do xantar acompañei o meu pai, que é asegurador, á casa da súa
clientela. Primeiro visitamos o \VUgras{ferreiro} \textsuperscript{(5)}, que vive
preto da pousada. Estaba ocupado ferrando un cabalo. Admirei a forza do seu compañeiro,
que mantiña firme sobre o seu mandil de coiro o \VUgras{casco} do cabalo, mentres o
ferreiro fixaba a \VUgras{ferradura} \textsuperscript{(6)} con fortes marteladas.

%%K t07-c1-06-1 p
6. --- Despois fomos á casa do \VUgras{zapateiro} \textsuperscript{(7)} da aldea, o
vello Xacobe. Aínda que ten por letreiro unha \VUgras{bota} fermosísima, contentábase
con botar solas novas a un vello par de zapatos furados.

%%K t07-c1-06-2 p
O vello díxonos rindo: «Na vila era “fabricante de zapatos” e non tiña clientes. Aquí
son “compoñedor de zapatos” e o traballo é moito.»

%%K t07-c1-07-1 p
7. --- Faloume do seu veciño o \VUgras{barbeiro} \textsuperscript{(8)}, coa clientela
do cal a xente non está contenta. As súas \VUgras{navallas} sempre están mocadas e as
súas \VUgras{tesoiras} non cortan ben en ningún momento.

Pero, como é o único barbeiro da aldea, un vese obrigado, claro está, a facerse
afeitar e cortar o pelo por el. Ademais é home avaro e desagradable.

%%K t07-c1-08-1 p
8. --- Por sorte, os demais \VUgras{veciños} non se lle parecen. Por exemplo o
\VUgras{carreteiro} \textsuperscript{(9)} é un home bondadoso, traballador e
servizal. É un pracer facerlle amañar un carro. O seu traballo sempre está ben
rematado.

%%K t07-c1-09-1 p
9. --- O vello Xacobe tampouco se queixou do \VUgras{albardeiro}
\textsuperscript{(10)}, ao que ás veces axuda a coser os \VUgras{arreos} cando urxe o
traballo.

%%K t07-c1-09-2 p
O meu pai preguntoulle pola súa familia: «Todos están ben. A miña neta está na
\VUgras{escola} \textsuperscript{(12)}. A súa \VUgras{mestra}
\textsuperscript{(13)} está moi contenta con ela; é boa \VUgras{alumna}
\textsuperscript{(14)}. Non se parece ao lercho do meu fillo; el só pensa en xogar. O
\VUgras{mestre} \textsuperscript{(15)} non consegue que aprenda nada.»

%%K t07-tit-4 sub
\VUsaut{3.0mm}
\VUpk{10.2pt}{40}{Falatorio da aldea.}
\VUsaut{3.0mm}

%%K t07-c1-10-1 p
10. --- O vello contounos despois as novas da aldea. Van construír unha escola nova,
pois a que está instalada na \VUgras{casa do concello} quedou pequena de máis.
Ademais é doado, pois abondará con mercar unha parte da horta do \VUgras{muiñeiro.}
Iso será moi proveitoso para o pobre home. Por mor do frío, as \VUgras{moas} do
\VUgras{muíño} \textsuperscript{(16)} xa non poden moer o trigo, pois a
\VUgras{roda} \textsuperscript{(17)} quedou parada polo \VUgras{xeo}
\textsuperscript{(25)} do \VUgras{regato} \textsuperscript{(23)}.

%%K t07-c1-11-1 p
11. --- «Falando de construcións, vostede viu a nosa nova \VUgras{igrexa}
\textsuperscript{(18)}? É moi pintoresca baixo o seu manto de neve. Os
\VUgras{corvos} \textsuperscript{(19)}, que xa non recoñecen o seu vello campanario,
pousan tristes na grande árbore do \VUgras{cemiterio} \textsuperscript{(20)}.»

%%K t07-c1-12-1 p
12. --- Esa idea do cemiterio fíxolle lembrar ao zapateiro as persoas que o meu pai
coñecera e que morreron o ano pasado. «Cantas \VUgras{sepulturas}
\textsuperscript{(21)}, meu bo señor, se engadiron ás outras, baixo o \VUgras{ciprés}
\textsuperscript{(22)}! A morte segou o noso vello notario. Todos os aldeáns asistiron
ao seu \VUgras{enterro.}»

%%K t07-c1-13-1 p
13. --- «Velaquí o seu sucesor, que patina no regato. Acaba de vir facerme amañar a
correa do seu \VUgras{patín} \textsuperscript{(26)}, que estaba rota.»

%%K t07-c1-14-1 p
14. --- «Velaquí tamén, na ribeira do regato, o novo \VUgras{garda rural}
\textsuperscript{(27)}. É un ex-gardacivil que aterra os pillabáns da aldea. Foxen
axiña que ven as súas \VUgras{polainas} \textsuperscript{(29)} e o seu \VUgras{quepis}
\textsuperscript{(28)}.»

%%K t07-c1-15-1 p
15. --- «Ha! Velaquí o vello Xosé, que leva un \VUgras{carriño de feixes de leña}
\textsuperscript{(30)} para o panadeiro. --- É verdade --- dixo meu pai ---;
recoñezo a súa xunta de \VUgras{mulas} \textsuperscript{(31)}. Ben preciso falar con
el; vou aproveitar a ocasión. Ata máis ver, tío Xacobe.»

%%K t07-c1-16-1 p
16. --- Xusto cando estabamos a saír, os nenos da aldea deixaron a escola e botáronse correndo
ao \VUgras{escorregadoiro} \textsuperscript{(33)} co risco de caeren e de romperen un
membro na \VUgras{caída} \textsuperscript{(34)}. Un deles colleu o seu
\VUgras{trineo} \textsuperscript{(32)} na casa do carreteiro, sentou nel un dos seus
compañeiros e marchou correndo.

%%K t07-c1-17-1 p
17. --- Outros botáronse a un \VUgras{montón de neve} \textsuperscript{(37)} preto da
\VUgras{ponte} \textsuperscript{(40)} e comezaron a bombardear a \VUgras{estatua de
neve} \textsuperscript{(39)} que ergueran o día antes e á que lle puxeran un sombreiro,
unha pipa e unha mochila en bandoleira.

%%K t07-c1-18-1 p
18. --- Pouco lles importa o vento xeado e o frío. A meirande parte nin sequera teñen
\VUgras{bufanda} \textsuperscript{(35)} e, para quentaren de novo despois da batalla
coas \VUgras{bólas de neve} \textsuperscript{(38)}, abóndalles unha presa de
\VUgras{castañas} \textsuperscript{(42)} moi quentes mercadas á vella
\VUgras{vendedora} Xacobina, que trema de frío baixo o seu \VUgras{mantelo}
\textsuperscript{(43)} delgado de máis.

%%K t07-tit-5 sub
\VUsaut{3.0mm}
\VUcentre{12.0pt}{0}{\textit{Segunda escena.}}
\VUsaut{3.0mm}

%%K t07-tit-6 sub
\VUpk{10.2pt}{40}{A casa labrega na primavera.}
\VUsaut{4.0mm}

%%K t07-c2-01-1 p
1. --- Malia todos os praceres do inverno, saudamos con fonda ledicia a volta da
\VUgras{primavera.}

%%K t07-c2-01-2 p
Que cadro tan alegre semella un patio de granxa, no serán, no mes de maio!

%%K t07-c2-02-1 p
2. --- No horizonte o \VUgras{sol} \textsuperscript{(1)} déitase amodo, asolagando o
\VUgras{campo} \textsuperscript{(3)} cos seus \VUgras{raios} purpúreos
\textsuperscript{(2)}. O dilixente \VUgras{labrador} \textsuperscript{(6)} apura
para arar a súa \VUgras{leira} \textsuperscript{(4)}.

%%K t07-c2-02-2 p
Co seu \VUgras{arado} \textsuperscript{(5)} enganchado a un vigoroso \VUgras{cabalo}
\textsuperscript{(7)}, traza o \VUgras{rego} \textsuperscript{(8)} no que o
\VUgras{sementador} \textsuperscript{(9)} botará a mans cheas a \VUgras{semente} que
dará as espigas douradas.

%%K t07-c2-03-1 p
3. --- Os \VUgras{segadores} \textsuperscript{(11)} provistos das súas fouces
segaron a herba do prado, os xuntadores secaron o \VUgras{feo}
\textsuperscript{(10)} dándolle voltas coas \VUgras{forcadas} \textsuperscript{(12)} e
cos angazos, e velaquí que volven.

%%K t07-c2-03-2 p
Uns camiñan diante do pesado carro tirado por bois; levan a súa ferramenta ao ombreiro.
Outros, máis lacazáns, instaláronse comodamente sobre o feo recendente.

%%K t07-c2-04-1 p
4. --- Ao pasaren fan un saúdo amigable ao \VUgras{xardineiro}
\textsuperscript{(16)} ocupado na \VUgras{horta das froitas} \textsuperscript{(13)}
en podar as \VUgras{árbores froiteiras} e enxertar as \VUgras{pereiras}
\textsuperscript{(14)}. As \VUgras{ameixeiras} \textsuperscript{(15)} comezan a
florecer e, na \VUgras{horta dos legumes} \textsuperscript{(17)} ben estercada, os
legumes, as coles e as ensaladas xa están medradas.

%%K t07-c2-04-2 p
No muriño, un belo \VUgras{pavón} \textsuperscript{(18)} desprega a súa cola para facer
admirar as súas plumas fermosísimas.

%%K t07-tit-7 sub
\VUpk{10.2pt}{40}{O patio da granxa.}
\VUsaut{3.0mm}

%%K t07-c2-05-1 p
5. --- As portas da \VUgras{cabaleriza} \textsuperscript{(19)} están abertas e vese o
comedeiro e a \VUgras{cama de palla} \textsuperscript{(23)} da \VUgras{egua}
\textsuperscript{(21)} que o \VUgras{cabaleirizo} \textsuperscript{(20)} acaba de
desenganchar. O \VUgras{poldro} \textsuperscript{(22)}, tan cómico coas súas pernas
longas, segue a súa nai dando chimpos.

%%K t07-c2-06-1 p
6. --- Preto do \VUgras{pozo} \textsuperscript{(29)} as \VUgras{vacas}
\textsuperscript{(25)} van beber na \VUgras{artesa} de madeira
\textsuperscript{(26)} que lles serve de bebedoiro. O \VUgras{xato}
\textsuperscript{(24)} aproveita a ocasión para mamar, e o \VUgras{boi}
\textsuperscript{(27)}, ao cheirar a herba da meda, brúa ledo e ergue con orgullo a
cabeza cos seus \VUgras{cornos} \textsuperscript{(28)} poderosos.

%%K t07-c2-07-1 p
7. --- Todos traballan. A \VUgras{criada} \textsuperscript{(32)} ten na man a
\VUgras{corda} \textsuperscript{(31)} do pozo. Tira auga cun \VUgras{balde}
\textsuperscript{(30)} para dar de beber ao gando. Preto da \VUgras{palleira}
\textsuperscript{(33)}, un home rastra as pallas, e diante da porta da \VUgras{casa}
\textsuperscript{(34)} unha vella \VUgras{fiandeira} \textsuperscript{(35)}, coa súa
roda de fiar e a súa \VUgras{roca}, fía o \VUgras{liño} ou o \VUgras{cánabo} como no
tempo de antano.

%%K t07-tit-8 sub
\VUsaut{3.0mm}
\VUpk{10.2pt}{40}{O granxeiro.}
\VUsaut{3.0mm}

%%K t07-c2-08-1 p
8. --- O \VUgras{granxeiro} \textsuperscript{(36)} dirixe todos eses traballos e dá
instrucións aos seus criados. Recomenda gardar baixo teito a \VUgras{grade}
\textsuperscript{(40)} e o \VUgras{sacho} \textsuperscript{(39)} que esqueceron preto
da \VUgras{caseta} \textsuperscript{(38)} do \VUgras{can} \textsuperscript{(37)}.

%%K t07-c2-09-1 p
9. --- Os seus berros asustan unha \VUgras{pomba} \textsuperscript{(41)}, que voa a
todo bater ao \VUgras{pombal} \textsuperscript{(42)}, e as \VUgras{galiñas}
\textsuperscript{(43)}, que apuran a volver ao \VUgras{galiñeiro}
\textsuperscript{(44)}.

%%K t07-c2-10-1 p
10. --- Diante do \VUgras{cortello das ovellas} \textsuperscript{(48)}, velaquí o
vello \VUgras{pastor} \textsuperscript{(45)} que trae de volta o seu \VUgras{rabaño}
\textsuperscript{(49)}. Coa súa \VUgras{cachaba} \textsuperscript{(46)}, que nunca
lle falta, o mesmo ca a súa \VUgras{blusa} \textsuperscript{(47)} descolorida, leva ao
curro os \VUgras{carneiros} \textsuperscript{(50)}, as \VUgras{ovellas}
\textsuperscript{(53)}, os \VUgras{años} \textsuperscript{(54)}, as \VUgras{cabras}
\textsuperscript{(51)} e os \VUgras{cabritos} \textsuperscript{(52)}. O seu fiel
\VUgras{can} \textsuperscript{(55)} Argos vén o derradeiro e apura cos seus ladridos a
dilixencia dos que se atrasan.

%%K t07-c2-11-1 p
11. --- Todo ese ruído e ese movemento non turban a quietude da boa \VUgras{vaca
leiteira} \textsuperscript{(56)}, que fai tintinar a súa \VUgras{campaíña}
\textsuperscript{(57)} mentres a muxen diante da porta do \VUgras{establo}
\textsuperscript{(58)}.

%%K t07-c2-12-1 p
12. Semella comprender que ten que dar o seu leite para que a \VUgras{granxeira}
\textsuperscript{(60)} poida facer \VUgras{manteiga} na \VUgras{manteigueira}
\textsuperscript{(61)} ou os excelentes \VUgras{queixos} \textsuperscript{(64)}, que
escorren da táboa pousada sobre a \VUgras{cuba} \textsuperscript{(63)}.

%%K t07-c2-13-1 p
13. --- Toda a \VUgras{leitaría} \textsuperscript{(59)} mantense moi limpa polo
\VUgras{criado da granxa} \textsuperscript{(65)}, que acaba de lavar as
\VUgras{leiteiras} \textsuperscript{(62)} no gran balde que leva na man.

%%K t07-tit-9 sub
\VUsaut{3.0mm}
\VUpk{10.2pt}{40}{O galiñeiro.}
\VUsaut{3.0mm}

%%K t07-c2-14-1 p
14. --- Cos seus grosos zocos, camiña algo pesadamente, e así fai fuxir o
\VUgras{pavo} \textsuperscript{(68)}, as \VUgras{patas} \textsuperscript{(66)}, os
\VUgras{patos} \textsuperscript{(67)} e os \VUgras{gansos} \textsuperscript{(69)}, que
abren de par en par \VUgras{o seu bico} \textsuperscript{(70)} e berran inquedos.

%%K t07-c2-14-2 p
O \VUgras{galo} \textsuperscript{(71)}, máis loitador, ergue a súa \VUgras{crista}
\textsuperscript{(72)} vermella, sacode as plumas da súa \VUgras{cola}
\textsuperscript{(73)} e fai oír un «quiquiriquí» sonoro.

%%K t07-c2-15-1 p
15. --- Unha \VUgras{galiña nai} \textsuperscript{(74)} anda a rapinar no
\VUgras{esterco} \textsuperscript{(81)} cos seus \VUgras{pitiños}
\textsuperscript{(75)}. O \VUgras{leitón} \textsuperscript{(78)} foxe cara á súa nai,
a enorme \VUgras{porca} \textsuperscript{(77)}, que vemos preto do cortello dos
porcos.

%%K t07-c2-16-1 p
16. --- Nas súas gaiolas, os \VUgras{coellos} \textsuperscript{(79)} comen con ansia a
\VUgras{herba} \textsuperscript{(80)} tenra que lles trae a filla do granxeiro.
Xoaniña quere moito os seus coellos e, cada día, despois de ir buscar os
\VUgras{ovos} \textsuperscript{(76)} frescos ao galiñeiro, vén visitar os seus
amiguiños.
\VUsaut{6.0mm}
\VUfilet{22mm}
